% generated by roam.py handbook export -- не редактировать, править заметки
% 20261004153518-монотонность_унимодальнои_функции.org (body 3fb799ddce5a)

\begin{theorem}[{Монотонность унимодальной функции}]\label[theorem]{note:1D154EFA-62D8-4BC3-B9CD-CCDE31F166B0}
  Пусть $\phi(t)$~--- \hyperref[note:CD52D818-AA25-471C-97C2-5DBE9787CE53]{унимодальна} на $[a;b]$ и $x^{*}$~--- точка минимума,
  тогда для любых $x_1 < x_2$ из $[a;b]$ верно:

  \begin{enumerate}
    \item $\phi(x_1) < \phi(x_2) \Rightarrow x^{*} \in (a;x_2]$
    \item $\phi(x_1) > \phi(x_2) \Rightarrow x^{*} \in [x_1; b)$
    \item $\phi(x_1) = \phi(x_2) \Rightarrow x^{*} \in (x_1; x_2)$
  \end{enumerate}
\end{theorem}

\begin{proof}
  Если $\phi(x_1) < \phi(x_2)$. Допустим $x^{*} \in (x_2;b)$, но тогда $x_1,x_2$ лежат на промежутке, где функция строго убывает. Противоречие. Аналогично для пункта 2.

  Для равенства. Тоже пойдем от обратного. Пусть $x^{*} \in (a; x_1]$, но тогда $x_1,x_2$ лежат на одном промежутке строгой монотонности и $\phi(x_1) \ne \phi(x_2)$. Аналогично для $[x_2;b)$.
\end{proof}
